% ============================================================
% Supplementary File S3: U-T-I-O Scoring Rubric
% ============================================================

\newpage
\section*{Supplementary File S3: U--T--I--O Scoring Rubric for Multi-Case Analysis}
\label{sec:S3}

\noindent\textbf{Purpose:} Provide explicit, reproducible criteria for assigning each case a
High~/~Medium~/~Low--Medium~/~Low rating on each of the four U-T-I-O domains.
Enables a second rater or replicator to reproduce or challenge the case ratings
reported in the manuscript.

\noindent\textbf{Scope:} Applied to four cases in this review (U.S.\ ONC HIE; NHS Shared Care
Records; Estonia eHealth; U.S.\ VA/DoD EHR modernization) but designed to generalise
to any national or organisational eHealth implementation.

\noindent\textbf{Rating scale:} \textbf{H}~=~High;\; \textbf{M}~=~Medium;\;
\textbf{L--M}~=~Low--Medium (constrained by clear gaps but not absent);\;
\textbf{L}~=~Low.

\bigskip

% ------------------------------------------------------------------
% U DOMAIN
% ------------------------------------------------------------------

\subsection*{U --- User Readiness}

\noindent The extent to which clinicians, administrators, and patients are prepared and
willing to engage with the eHealth system.

\renewcommand{\arraystretch}{1.3}
\setlength{\tabcolsep}{4pt}

\begin{table}[H]
\footnotesize
\begin{tabular}{p{2.5cm} p{3.1cm} p{3.1cm} p{3.1cm} p{3.1cm}}
\toprule
\rowcolor{headerblue}
\textcolor{white}{\textbf{Construct}} &
\textcolor{white}{\textbf{High (H)}} &
\textcolor{white}{\textbf{Medium (M)}} &
\textcolor{white}{\textbf{Low--Medium (L--M)}} &
\textcolor{white}{\textbf{Low (L)}} \\
\midrule

\rowcolor{rowgray}
Usability &
Documented favourable usability scores; SUS/clinician-satisfaction surveys consistently $\geq$75 &
Mixed survey results; some pain points but no widespread crisis &
Clinician satisfaction below targets in published audits &
Persistent, documented usability crisis (e.g., GAO findings of dissatisfaction) \\

Workflow integration &
Designed with clinician co-design; embedded in existing workflows &
Partial workflow embedding; some retrofit required &
Workflow disruption acknowledged in audits &
Documented workflow disruption with productivity loss \\

\rowcolor{rowgray}
Training adequacy &
Funded, multi-modal training programme deployed pre-go-live &
Training exists but variable quality across sites &
Training gaps identified in audits &
No structured training programme \\

Trust / perceived usefulness &
Clinician advocates speak publicly in favour of the system &
Mixed clinician advocacy &
Clinician resistance documented in formal surveys &
Active clinician opposition or resistance \\

\rowcolor{rowgray}
Patient engagement &
Patient portal adoption $\geq$50\% nationally &
Patient portal exists; adoption variable &
Portal exists but low usage &
No patient-facing functionality \\

\bottomrule
\end{tabular}
\end{table}

\noindent\textbf{Rating decision rule:}
Two or more constructs meeting H criteria with no constructs at L $\rightarrow$ \textbf{H}.
Mix of M with one or two H/L $\rightarrow$ \textbf{M}.
Two or more constructs at L--M with usability or training gaps documented in audits $\rightarrow$ \textbf{L--M}.
Documented persistent user-satisfaction crisis (e.g., GAO findings) $\rightarrow$ \textbf{L}.

\bigskip

% ------------------------------------------------------------------
% T DOMAIN
% ------------------------------------------------------------------

\subsection*{T --- Technical Interoperability}

\noindent The capacity of eHealth systems to exchange, interpret, and use clinical data
across organisational boundaries.

\begin{table}[H]
\footnotesize
\begin{tabular}{p{2.5cm} p{3.1cm} p{3.1cm} p{3.1cm} p{3.1cm}}
\toprule
\rowcolor{headerblue}
\textcolor{white}{\textbf{Construct}} &
\textcolor{white}{\textbf{High (H)}} &
\textcolor{white}{\textbf{Medium (M)}} &
\textcolor{white}{\textbf{Low--Medium (L--M)}} &
\textcolor{white}{\textbf{Low (L)}} \\
\midrule

\rowcolor{rowgray}
Semantic interoperability &
National terminology binding (SNOMED CT, LOINC); enforced &
Standards adopted but not nationally bound &
Partial standards adoption &
No common terminology \\

Syntactic interoperability &
HL7 FHIR R4 + CDA mandated; conformance testing &
FHIR adoption variable across vendors &
Legacy formats predominant &
No interoperability standard \\

\rowcolor{rowgray}
System quality &
Documented uptime $\geq$99.5\%; mature tooling &
Acceptable uptime; episodic incidents &
Reliability concerns in audits &
System reliability crisis \\

Data completeness &
National data set obligation enforced &
Voluntary upload with high participation &
Variable coverage &
Sparse coverage \\

\rowcolor{rowgray}
API / exchange layer &
National exchange backbone deployed (e.g., X-Road) &
Federated exchange (HIE network) operational &
Bilateral interfaces only &
No exchange \\

Integration burden &
Vendor agnostic; standards-driven integration &
Some vendor lock-in but multiple options &
Heavy custom integration required &
Vendor monoculture / lock-in \\

\bottomrule
\end{tabular}
\end{table}

\noindent\textbf{Rating decision rule:}
National exchange backbone + enforced standards + uptime $\rightarrow$ \textbf{H}.
Federated/standards-based but voluntary $\rightarrow$ \textbf{M}.
Standards adopted but with substantial integration debt $\rightarrow$ \textbf{L--M}.
No common standards / pervasive bilateral interfaces $\rightarrow$ \textbf{L}.

\bigskip

% ------------------------------------------------------------------
% I DOMAIN
% ------------------------------------------------------------------

\subsection*{I --- Institutional Alignment}

\noindent The degree to which regulatory, policy, governance, and legal environments
support and incentivise eHealth adoption.

\begin{table}[H]
\footnotesize
\begin{tabular}{p{2.5cm} p{3.1cm} p{3.1cm} p{3.1cm} p{3.1cm}}
\toprule
\rowcolor{headerblue}
\textcolor{white}{\textbf{Construct}} &
\textcolor{white}{\textbf{High (H)}} &
\textcolor{white}{\textbf{Medium (M)}} &
\textcolor{white}{\textbf{Low--Medium (L--M)}} &
\textcolor{white}{\textbf{Low (L)}} \\
\midrule

\rowcolor{rowgray}
Regulatory mandate &
Statutory obligation to participate (e.g., national EHR law) &
Strong incentive programme (HITECH-style) without mandate &
Voluntary policy framework &
No coordinating policy \\

National digital health strategy &
Documented, funded, multi-year national strategy &
Policy framework with episodic funding &
Strategy in development &
No national strategy \\

\rowcolor{rowgray}
Information governance &
Mature data-sharing agreements; national opt-out &
Information governance frameworks exist; local variation &
Information governance gaps acknowledged &
No coherent governance \\

Privacy law alignment &
Health-IT-specific privacy law (HIPAA, GDPR + national act) &
Generic privacy law applied to health data &
Privacy gaps acknowledged &
No applicable framework \\

\rowcolor{rowgray}
Interoperability mandate &
Information blocking prohibited and enforced &
Interoperability standards required (e.g., Cures Act) &
Voluntary interoperability targets &
No interoperability obligation \\

Incentive alignment &
Direct financial incentives for adoption + use &
Indirect incentives (quality reporting) &
Limited incentive structure &
No incentive structure \\

\bottomrule
\end{tabular}
\end{table}

\noindent\textbf{Rating decision rule:}
Statutory mandate + enforced governance + national strategy $\rightarrow$ \textbf{H}.
Strong incentives + interoperability rules without mandate $\rightarrow$ \textbf{M} to \textbf{H} depending on enforcement.
Voluntary frameworks with documented gaps $\rightarrow$ \textbf{L--M}.
Absent or fragmented regulatory environment $\rightarrow$ \textbf{L}.

\bigskip

% ------------------------------------------------------------------
% O DOMAIN
% ------------------------------------------------------------------

\subsection*{O --- Organizational Execution}

\noindent The internal organisational capabilities and leadership practices that determine whether an
eHealth adoption strategy is planned, resourced, governed, and executed effectively.

\begin{table}[H]
\footnotesize
\begin{tabular}{p{2.5cm} p{3.1cm} p{3.1cm} p{3.1cm} p{3.1cm}}
\toprule
\rowcolor{headerblue}
\textcolor{white}{\textbf{Construct}} &
\textcolor{white}{\textbf{High (H)}} &
\textcolor{white}{\textbf{Medium (M)}} &
\textcolor{white}{\textbf{Low--Medium (L--M)}} &
\textcolor{white}{\textbf{Low (L)}} \\
\midrule

\rowcolor{rowgray}
Leadership commitment &
Sustained executive sponsorship; named accountable owners &
Leadership engagement at launch; sustainment variable &
Leadership turnover or disengagement during execution &
No accountable owner \\

Implementation planning &
Documented multi-phase plan with milestones; published &
Plan exists; milestones missed without remediation &
Reactive rather than planned execution &
No formal plan \\

\rowcolor{rowgray}
Change management &
Funded change-management capability throughout lifecycle &
Change-management activities at go-live only &
Change management identified as gap in audits &
No change-management capacity \\

Stakeholder engagement &
Clinician co-design + ongoing feedback channels &
Stakeholder engagement at launch; episodic after &
Stakeholder voice limited in execution &
No engagement mechanism \\

\rowcolor{rowgray}
Vendor governance &
Mature vendor-management practice; contract performance tracked &
Contract management functional &
Vendor disputes documented in audits &
Vendor relationship dysfunctional \\

Resource capacity &
IT staff, financial, and clinical-time resources adequate &
Resource constraints acknowledged but managed &
Resource gaps cited as adoption barrier &
Severe resource constraint \\

\rowcolor{rowgray}
Issue management &
Proactive issue identification and resolution &
Reactive issue resolution &
Issue backlog documented in audits &
No issue-management process \\

\bottomrule
\end{tabular}
\end{table}

\noindent\textbf{Rating decision rule:}
All seven constructs meeting at least M, with leadership and change-management at H $\rightarrow$ \textbf{H}.
Mixed M/L--M with vendor or change-management weakness $\rightarrow$ \textbf{M}.
Two or more constructs at L--M, audit findings of execution gaps $\rightarrow$ \textbf{L--M}.
Documented severe execution failures (cost overruns, schedule slips, audit findings of inadequate management) $\rightarrow$ \textbf{L}.

\bigskip

% ------------------------------------------------------------------
% APPLICATION TO FOUR CASES
% ------------------------------------------------------------------

\subsection*{Evidence-Anchored Ratings for the Four Cases}

\noindent For each case, ratings are anchored to specific sources from the bibliography
(citation keys in parentheses). This is the audit trail a reviewer or replicator can
use to challenge or reproduce any rating.

\medskip

\noindent\textbf{Case 1: U.S.\ ONC HIE --- U(M) T(H) I(H) O(M)}

\begin{table}[H]
\footnotesize
\begin{tabular}{p{1.2cm} >{\centering\arraybackslash}p{1.4cm} p{12.5cm}}
\toprule
\rowcolor{headerblue}
\textcolor{white}{\textbf{Domain}} &
\textcolor{white}{\textbf{Rating}} &
\textcolor{white}{\textbf{Evidence anchor}} \\
\midrule
\rowcolor{rowgray}
U & \cellcolor{medamber}\textcolor{medambertext}{\textbf{M}} &
\cite{onc2025reports} documents variable clinician engagement and adoption-rate gaps across participating entities; \cite{eden2016barriers} identifies clinician trust and workflow concerns as persistent HIE barriers; \cite{kruse2016adoption} lists usability and training as adoption inhibitors. \\
T & \cellcolor{highgreen}\textcolor{highgreentext}{\textbf{H}} &
\cite{onc2025reports} references FHIR API mandates under the Cures Act; \cite{holmgren2023policyhie} describes the U.S.\ as one of five mature HIE-policy nations; nationwide TEFCA infrastructure deployed. \\
\rowcolor{rowgray}
I & \cellcolor{highgreen}\textcolor{highgreentext}{\textbf{H}} &
HITECH Act incentive program drove U.S.\ EHR adoption from $\sim$10\% to $>$80\% (\cite{onc2025reports}); 21st Century Cures Act information-blocking rules now enforced (\cite{holmgren2023policyhie}); strong statutory framework. \\
O & \cellcolor{medamber}\textcolor{medambertext}{\textbf{M}} &
\cite{onc2025reports} documents variable participation rates across HIEs; organisational uptake heterogeneity is the primary O-domain weakness; \cite{fennelly2020national} describes execution as the differentiator across national programs. \\
\bottomrule
\end{tabular}
\end{table}

\medskip

\noindent\textbf{Case 2: NHS Shared Care Records (UK) --- U(M) T(H) I(H) O(M)}

\begin{table}[H]
\footnotesize
\begin{tabular}{p{1.2cm} >{\centering\arraybackslash}p{1.4cm} p{12.5cm}}
\toprule
\rowcolor{headerblue}
\textcolor{white}{\textbf{Domain}} &
\textcolor{white}{\textbf{Rating}} &
\textcolor{white}{\textbf{Evidence anchor}} \\
\midrule
\rowcolor{rowgray}
U & \cellcolor{medamber}\textcolor{medambertext}{\textbf{M}} &
\cite{nhs2025sharedcare} acknowledges variable local clinician engagement; sustained training and workflow integration cited as central challenges. \\
T & \cellcolor{highgreen}\textcolor{highgreentext}{\textbf{H}} &
\cite{nhs2025sharedcare} describes national interoperability standards plus a shared-care-record platform; cross-organisational information sharing operational. \\
\rowcolor{rowgray}
I & \cellcolor{highgreen}\textcolor{highgreentext}{\textbf{H}} &
NHS Long Term Plan + national data opt-out + information governance frameworks documented in \cite{nhs2025sharedcare}; central national strategy authority. \\
O & \cellcolor{medamber}\textcolor{medambertext}{\textbf{M}} &
\cite{nhs2025sharedcare} flags local information-governance variability and workflow integration as persistent execution challenges; \cite{fennelly2020national} lists O-domain factors as decisive across national EHR rollouts. \\
\bottomrule
\end{tabular}
\end{table}

\medskip

\noindent\textbf{Case 3: Estonia eHealth --- U(H) T(H) I(H) O(H)}

\begin{table}[H]
\footnotesize
\begin{tabular}{p{1.2cm} >{\centering\arraybackslash}p{1.4cm} p{12.5cm}}
\toprule
\rowcolor{headerblue}
\textcolor{white}{\textbf{Domain}} &
\textcolor{white}{\textbf{Rating}} &
\textcolor{white}{\textbf{Evidence anchor}} \\
\midrule
\rowcolor{rowgray}
U & \cellcolor{highgreen}\textcolor{highgreentext}{\textbf{H}} &
\cite{estonia2026ehealth} reports near-comprehensive coverage and high clinician engagement; 99\% patient e-Health Record adoption widely cited. \\
T & \cellcolor{highgreen}\textcolor{highgreentext}{\textbf{H}} &
\cite{estonia2026ehealth} describes the X-Road secure data-exchange layer; \cite{holmgren2023policyhie} cites Estonia as architectural reference; enforced national standards. \\
\rowcolor{rowgray}
I & \cellcolor{highgreen}\textcolor{highgreentext}{\textbf{H}} &
\cite{europeancommission2016estoniaehr} documents the legal mandate for provider participation; 2008 Estonian Health Information System Act provides statutory backbone. \\
O & \cellcolor{highgreen}\textcolor{highgreentext}{\textbf{H}} &
\cite{estonia2026ehealth} describes sustained ministry-level governance and evolving operational frameworks; consistent execution across two decades of national infrastructure operation. \\
\bottomrule
\end{tabular}
\end{table}

\medskip

\noindent\textbf{Case 4: U.S.\ VA/DoD EHR Modernization --- U(L--M) T(L--M) I(H) O(L--M)}

\begin{table}[H]
\footnotesize
\begin{tabular}{p{1.2cm} >{\centering\arraybackslash}p{1.4cm} p{12.5cm}}
\toprule
\rowcolor{headerblue}
\textcolor{white}{\textbf{Domain}} &
\textcolor{white}{\textbf{Rating}} &
\textcolor{white}{\textbf{Evidence anchor}} \\
\midrule
\rowcolor{rowgray}
U & \cellcolor{lowred}\textcolor{lowredtext}{\textbf{L--M}} &
\cite{gao2025vaehr} documents persistent clinician dissatisfaction at deployed VA sites; \cite{gao2024dodehr} cites user-satisfaction scores below targets at DoD sites. The rubric's U decision rule lists ``Documented persistent user-satisfaction crisis (e.g., GAO findings)~$\rightarrow$~L''; the L--M rating reflects the documented crisis tempered by ongoing remediation activity at deployed sites. \textit{Revised from M (2026-04-30); see audit-trail note below.} \\
T & \cellcolor{lowred}\textcolor{lowredtext}{\textbf{L--M}} &
\cite{gao2025vaehr} describes mature commercial Oracle Health platform but substantial configuration burden integrating with legacy VistA and DoD clinical systems; documented system-quality and data-completeness gaps at deployed sites. The rubric's T decision rule ``Standards adopted but with substantial integration debt~$\rightarrow$~L--M'' is the closest match. \textit{Revised from M (2026-04-30); see audit-trail note below.} \\
\rowcolor{rowgray}
I & \cellcolor{highgreen}\textcolor{highgreentext}{\textbf{H}} &
Congressional mandate via FY2018 NDAA; sustained appropriations; clear executive-branch authority --- \cite{gao2025vaehr}, \cite{gao2024dodehr} both confirm institutional backing. \\
O & \cellcolor{lowred}\textcolor{lowredtext}{\textbf{L--M}} &
\cite{gao2025vaehr} flags inadequate vendor governance and slow remediation; \cite{gao2024dodehr} cites issue-management process inadequacies and schedule slips. Persistent execution gaps drive the lowest-tier ratings in the four-case set. \\
\bottomrule
\end{tabular}
\end{table}

\smallskip
\noindent\textbf{Audit-trail note (2026-04-30):} The U and T ratings for VA/DoD were revised from M to L--M following the AI-assisted second-rater pass documented in Supplementary File~S4. The second rater applied the rubric's explicit decision-rule trigger language for both domains; the original ratings (M for both) were a band more lenient than the rubric strictly indicated. Revised ratings retain a middle position (L--M rather than the second rater's L for U) to reflect ongoing remediation activity at deployed sites. Inter-rater agreement statistics in Supplementary File~S4 are computed on the \textit{pre-revision} ratings, as is methodologically appropriate.

\medskip
\noindent\textbf{How to reproduce:} A second rater should retrieve the cited source(s) for each cell, apply the rubric criteria above, and assign a rating. Rating disagreements should be reported as an inter-rater reliability outcome.

\bigskip

% ------------------------------------------------------------------
% LIMITATIONS OF THE RUBRIC
% ------------------------------------------------------------------

\subsection*{Limitations of this Rubric}

\begin{enumerate}
\item \textbf{Single-rater application.} Ratings in the current manuscript were assigned by one reviewer (the author). Independent dual-rating with inter-rater reliability is not yet performed.
\item \textbf{Categorical resolution.} Four bands (H / M / L--M / L) are coarse. A reviewer wanting finer resolution could move to a 5- or 7-point ordinal scale.
\item \textbf{Construct weighting.} Each domain treats its constructs as equally informative; in practice, some constructs (e.g., regulatory mandate within~I) may dominate the rating. This is not formalised here.
\item \textbf{Evidence weighting.} Government audit findings and peer-reviewed literature are weighted equally; some reviewers might down-weight grey literature.
\end{enumerate}

\noindent\textit{Last updated: 2026-04-29}
